\section{ابرتفکیک بدون نمونهٔ آموزشی با یادگیری عمیق داخلی}
\label{sec:zssr}

\subsection{مقدمه و انگیزش}
\label{subsec:zssr_intro}

یادگیری عمیق در سال‌های اخیر جهشی چشمگیر در عملکرد ابرتفکیک تصویر
\LTRfootnote{Super-Resolution (SR)}
ایجاد کرده است \cite{dong2014srcnn, kim2016vdsr, lim2017edsr, ledig2017srgan}.
با این حال، روش‌های مبتنی بر یادگیری نظارت‌شده با محدودیتی اساسی
مواجه‌اند: آن‌ها فرض می‌کنند فرایند کاهش وضوح از تصاویر با کیفیت بالا
\LTRfootnote{High-Resolution (HR)}
به تصاویر کم‌کیفیت
\LTRfootnote{Low-Resolution (LR)}
از پیش مشخص و ثابت است. به‌طور معمول، هستهٔ
کاهش‌نمونه‌برداری یک هستهٔ دومکعبی
\LTRfootnote{Bicubic kernel}
با ضد‌اعوجاج است (دستور پیش‌فرض \texttt{imresize} نرم‌افزار متلب)
و فرض می‌شود هیچ‌گونه آرتیفکت اضافی مانند نویز حسگر، فشرده‌سازی تصویر
یا تابع انتشار نقطه‌ای\LTRfootnote{Point Spread Function (PSF)} غیرایده‌آل
وجود ندارد \cite{shocher2018zssr}.

در عمل، تصاویر کم‌کیفیت واقعی به‌ندرت از این فرضیات پیروی می‌کنند.
تصاویر دانلود‌شده از اینترنت، عکس‌های گرفته‌شده با تلفن همراه، تصاویر
تاریخی قدیمی و داده‌های تصویری زیستی هر یک فرایند تصویربرداری منحصر
به‌فردی دارند. در چنین شرایطی، روش‌های ابرتفکیک پیشرفته مبتنی بر
یادگیری نظارت‌شده عملکردی ضعیف از خود نشان می‌دهند. به‌عنوان مثال،
شبکهٔ عمیق EDSR \cite{lim2017edsr} که بر داده‌های ایده‌آل آموزش دیده،
هنگام مواجهه با تصاویری که هستهٔ کاهش‌نمونه‌برداری آن‌ها غیردومکعبی
است یا آرتیفکت‌های حاصل از نویز و فشرده‌سازی دارند، نتایجی تار و
بدون جزئیات تولید می‌کند.

این مسئله انگیزهٔ اصلی شوخر و همکاران \cite{shocher2018zssr} را برای
ارائهٔ روش \lr{ZSSR}\LTRfootnote{Zero-Shot Super-Resolution} فراهم ساخت.
ایدهٔ کلیدی این پژوهش بهره‌گیری از قدرت یادگیری عمیق است، بدون آنکه
به هیچ دادهٔ آموزشی خارجی یا پیش‌آموزش وابسته باشد. به جای آن، از
تکرارپذیری درونی اطلاعات در یک تصویر واحد استفاده می‌شود تا یک شبکهٔ
عصبی کانولوشنی خاص هر تصویر در زمان آزمون آموزش داده شود. تا جایی که
نویسندگان اطلاع دارند، ZSSR نخستین روش ابرتفکیک مبتنی بر CNN بدون
نظارت است.

\subsection{قدرت پیش‌بینی آمار درونی تصویر}
\label{subsec:internal_stats}

پایهٔ نظری روش ZSSR بر مشاهده‌ای بنیادین استوار است: تصاویر طبیعی
تکرارپذیری درونی قوی در داده‌ها دارند. پچ‌های کوچک تصویر (مثلاً
$5\times5$ یا $7\times7$) بارها و بارها در درون یک تصویر واحد تکرار
می‌شوند، هم در مقیاس یکسان و هم در مقیاس‌های مختلف
\cite{glasner2009sr, zontak2011internal}.

این مشاهده به‌صورت تجربی توسط گلاسنر و همکاران
\cite{glasner2009sr} و زونتاک و ایرانی \cite{zontak2011internal}
بر روی صدها تصویر طبیعی تأیید شده و نشان داده شده که تقریباً برای
هر پچ کوچک و در تقریباً هر تصویر طبیعی صادق است. این ویژگی پایه‌ای
بسیاری از روش‌های بدون نظارت برای بهبود تصویر از جمله ابرتفکیک بدون
نظارت \cite{glasner2009sr, freedman2011selfexamples, huang2015selfex}،
ابرتفکیک کور \cite{michaeli2013blind} و محو‌زدایی کور
\cite{michaeli2014deblur, bahat2017deblur} بوده است.

شوخر و همکاران به یکی از مهم‌ترین شواهد این قدرت پیش‌بینی درونی اشاره
می‌کنند: در تصویری از ساختمانی با بالکن‌های متعدد، نرده‌های بسیار
ریز بالکن‌های کوچک تنها از درون خود تصویر قابل بازسازی هستند، زیرا
شواهدی از وجود آن‌ها در بالکن‌های بزرگ‌تر (در مکان و مقیاس دیگر درون
همان تصویر) یافت می‌شود. این شواهد در هیچ مجموعه‌دادهٔ خارجی، هر چقدر
هم بزرگ باشد، وجود ندارد. روش‌های ابرتفکیک پیشرفته‌ای مانند
EDSR \cite{lim2017edsr} که بر داده‌های خارجی آموزش دیده‌اند، در
بازسازی این اطلاعات خاصِ تصویر ناموفق‌اند.

علاوه بر این، زونتاک و ایرانی \cite{zontak2011internal} به‌صورت
تجربی نشان دادند که آنتروپی درونی پچ‌ها در یک تصویر واحد
به‌مراتب کمتر از آنتروپی خارجی پچ‌ها در مجموعه‌ای عمومی از
تصاویر طبیعی است. این یافته بدان معناست که آمار درونی تصویر
اغلب قدرت پیش‌بینی قوی‌تری نسبت به آمار خارجی فراهم می‌کند.
این برتری به‌ویژه در شرایط عدم قطعیت بیشتر و تخریب‌های تصویری
شدیدتر تقویت می‌شود \cite{zontak2011internal, mosseri2013combining}.

با وجود مزایای روش‌های بدون نظارت مبتنی بر پچ، آن‌ها با محدودیت‌هایی
مواجه‌اند: معمولاً بر شباهت اقلیدسی ساده میان پچ‌های کوچک با
اندازهٔ ثابت (مثلاً $5\times5$) و جستجوی $K$-نزدیک‌ترین همسایه
تکیه می‌کنند. در نتیجه، توانایی تعمیم‌پذیری آن‌ها برای پچ‌هایی که در
تصویر LR وجود ندارند محدود است و نمی‌توانند معیارهای شباهت ضمنی
یادگیری‌شده یا اندازه‌های غیریکنواخت ساختارهای تکرارشونده در تصویر
را مدیریت کنند.

\subsection{شبکهٔ عصبی کانولوشنی خاص تصویر}
\label{subsec:image_specific_cnn}

روش ZSSR قدرت پیش‌بینی و آنتروپی پایین اطلاعات درونی خاص تصویر را با
قابلیت‌های تعمیم‌پذیری یادگیری عمیق ترکیب می‌کند. فرض کنید تصویر
آزمون $I$ داده شده و هیچ نمونهٔ خارجی برای آموزش در دسترس نیست.
در این حالت، یک شبکهٔ CNN خاص تصویر ساخته می‌شود که منحصراً
برای حل مسئلهٔ ابرتفکیک همین تصویر طراحی شده است.

\subsubsection{استخراج نمونه‌های آموزشی از تصویر}

ایدهٔ کلیدی آن است که شبکه بر روی نمونه‌هایی آموزش ببیند که صرفاً
از خود تصویر آزمون استخراج شده‌اند. این نمونه‌ها به‌صورت زیر تولید
می‌شوند: تصویر LR ورودی $I$ با ضریب مقیاس مطلوب $s$ کاهش‌نمونه‌برداری
می‌شود تا نسخه‌ای با وضوح پایین‌تر، یعنی $I{\downarrow s}$، به دست آید.
سپس شبکه آموزش می‌بیند تا تصویر $I$ را از نسخهٔ کاهش‌یافته‌اش
$I{\downarrow s}$ بازسازی کند. پس از آموزش، همان شبکه بر روی تصویر
ورودی $I$ اعمال می‌شود و خروجی با وضوح بالا، یعنی $I{\uparrow s}$،
تولید می‌گردد. از آنجا که شبکهٔ آموزش‌دیده کاملاً کانولوشنی است،
قادر به اعمال بر تصاویر با اندازه‌های مختلف خواهد بود.

\subsubsection{افزایش داده}

از آنجا که مجموعهٔ آموزشی تنها شامل یک نمونه (تصویر آزمون) است، از
روش‌های افزایش داده برای استخراج جفت‌نمونه‌های بیشتر استفاده می‌شود.
تصویر آزمون $I$ به نسخه‌های کوچک‌تر متعددی کاهش‌نمونه‌برداری می‌شود:
$I = I_0, I_1, I_2, \ldots, I_n$. این نسخه‌ها نقش
\emph{پدرهای HR}\LTRfootnote{HR fathers} را ایفا می‌کنند. سپس هر
پدر HR با ضریب مقیاس $s$ کاهش‌نمونه‌برداری می‌شود تا
\emph{فرزندان LR}\LTRfootnote{LR sons} ایجاد شوند. مجموعهٔ حاصل شامل
جفت‌نمونه‌های LR-HR خاص تصویر است.

علاوه بر کاهش مقیاس، هر جفت LR-HR با استفاده از ۴ دوران
($0^\circ, 90^\circ, 180^\circ, 270^\circ$) و بازتاب آینه‌ای در
جهت‌های افقی و عمودی غنی‌سازی می‌شود. این عملیات $8\times$ نمونهٔ
آموزشی خاص تصویر بیشتر اضافه می‌کند.

\subsubsection{ابرتفکیک تدریجی}

به‌منظور استحکام بیشتر و امکان اعمال ضرایب مقیاس بزرگ حتی بر تصاویر
LR بسیار کوچک، ابرتفکیک به‌صورت تدریجی انجام می‌شود
\cite{glasner2009sr, timofte2016seven}. الگوریتم برای چندین ضریب مقیاس
میانی $(s_1, s_2, \ldots, s_m = s)$ اجرا می‌شود. در هر مقیاس میانی
$s_i$، تصویر SR تولیدشده $\text{HR}_i$ و نسخه‌های
کاهش‌یافته/دوران‌یافتهٔ آن به‌عنوان پدرهای HR جدید به مجموعهٔ
آموزشی (که به‌تدریج در حال رشد است) افزوده می‌شوند. سپس این نمونه‌ها
با ضریب مقیاس بعدی $s_{i+1}$ کاهش‌نمونه‌برداری شده و جفت‌نمونه‌های
آموزشی LR-HR جدید تولید می‌شوند. این فرایند تا رسیدن به افزایش
وضوح کامل مطلوب $s$ ادامه می‌یابد.

\subsection{معماری و بهینه‌سازی شبکه}
\label{subsec:arch_optim}

شبکه‌های CNN نظارت‌شده که بر مجموعه‌ای بزرگ و متنوع از نمونه‌های
خارجی LR-HR آموزش می‌بینند، باید تنوع عظیم تمام روابط ممکن LR-HR
را در وزن‌های یادگیری‌شدهٔ خود رمزگذاری کنند. به همین دلیل، این
شبکه‌ها معمولاً بسیار عمیق و پیچیده هستند. در مقابل، تنوع روابط
LR-HR در یک تصویر واحد به‌مراتب کمتر است و می‌تواند توسط شبکه‌ای
بسیار کوچک‌تر و ساده‌تر رمزگذاری شود.

معماری پیشنهادی ZSSR یک شبکهٔ کاملاً کانولوشنی ساده با مشخصات زیر است:

\begin{itemize}
\item \textbf{عمق:} ۸ لایهٔ پنهان
\item \textbf{کانال‌ها:} هر لایه ۶۴ کانال
\item \textbf{تابع فعال‌سازی:} \lr{ReLU} بر هر لایه
\item \textbf{ورودی:} تصویر LR درون‌یابی‌شده به اندازهٔ خروجی
\item \textbf{یادگیری:} تنها باقی‌ماندهٔ\LTRfootnote{Residual} بین
تصویر LR درون‌یابی‌شده و پدر HR آن یادگیری می‌شود (مشابه
\cite{dong2014srcnn, kim2016vdsr})
\end{itemize}

برای بهینه‌سازی شبکه، از تابع هزینهٔ $L_1$ به‌همراه بهینه‌ساز
\lr{Adam} \cite{kingma2014adam} استفاده شده است. نرخ یادگیری اولیه
$0.001$ در نظر گرفته می‌شود. به‌صورت دوره‌ای، یک برازش خطی بر خطای
بازسازی انجام می‌شود؛ اگر انحراف استاندارد از شیب برازش خطی بیشتر
باشد، نرخ یادگیری بر ۱۰ تقسیم می‌شود. آموزش زمانی متوقف می‌شود که
نرخ یادگیری به $10^{-6}$ برسد. این استراتژی کاهش نرخ یادگیری را
می‌توان به‌صورت زیر خلاصه کرد:
\begin{equation}
\label{eq:lr_schedule}
\eta_{t+1} =
\begin{cases}
\eta_t / 10, & \text{اگر}\; \sigma_{\text{err}} > c \cdot |\text{slope}| \\
\eta_t, & \text{در غیر این صورت}
\end{cases}
\end{equation}
که $\eta_t$ نرخ یادگیری در دورهٔ $t$، $\sigma_{\text{err}}$ انحراف
استاندارد خطای بازسازی و $c$ یک ثابت آستانه‌ای است.

برای شتاب‌بخشیدن به مرحلهٔ آموزش و استقلال زمان اجرا از اندازهٔ
تصویر آزمون $I$، در هر تکرار یک برش تصادفی با اندازهٔ ثابت (معمولاً
$128 \times 128$) از یک جفت پدر-فرزند انتخابی تصادفی گرفته می‌شود
(مگر آنکه جفت تصویر انتخاب‌شده کوچک‌تر باشد). احتمال نمونه‌برداری
هر جفت LR-HR در هر تکرار آموزشی غیریکنواخت و متناسب با اندازهٔ
پدر HR است: هرچه نسبت اندازهٔ پدر HR به تصویر آزمون $I$ به ۱
نزدیک‌تر باشد، احتمال نمونه‌برداری آن بیشتر است. این طراحی بازتابی
از اعتبار بیشتر نمونه‌های HR غیرمصنوعی نسبت به نمونه‌های سنتز‌شده
است.

در نهایت، از تکنیک \emph{خودتجمیع هندسی}
\LTRfootnote{Geometric self-ensemble} مشابه \cite{lim2017edsr}
استفاده می‌شود. ۸ خروجی مختلف برای ۸ دوران و بازتاب تصویر آزمون $I$
تولید و سپس ترکیب می‌شوند. در ZSSR به جای میانگین از \emph{میانه}
استفاده می‌شود. این خروجی همچنین با تکنیک بازافکنی
\LTRfootnote{Back-projection} \cite{irani1991backproj, glasner2009sr}
ترکیب می‌شود: هر ۸ تصویر خروجی چند تکرار بازافکنی را طی کرده
و نهایتاً تصویر میانه نیز با بازافکنی اصلاح می‌گردد.

\paragraph{زمان اجرا.}
هرچند آموزش در زمان آزمون انجام می‌شود، میانگین زمان اجرا برای هر
تصویر ۵۴ ثانیه برای یک افزایش مقیاس واحد (بر GPU مدل \lr{Tesla K-80})
است. این زمان مستقل از اندازهٔ تصویر و ضریب مقیاس نسبی $s$ است (به
دلیل برش‌های تصادفی با اندازهٔ ثابت). با این حال، نتایج بهتر با
افزایش تدریجی وضوح به دست می‌آید. برای مثال، استفاده از ۶ ضریب مقیاس
میانی PSNR را حدوداً $0.2$~dB بهبود می‌بخشد، اما زمان اجرا را به حدود
۵ دقیقه افزایش می‌دهد. مقایسه با EDSR+ \cite{lim2017edsr} نشان
می‌دهد که بر همان بستر سخت‌افزاری، EDSR+ حدود ۲۰ ثانیه برای ابرتفکیک
$\times 2$ یک تصویر $200 \times 200$ نیاز دارد، اما زمان اجرای آن
به‌صورت مربعی با اندازهٔ تصویر رشد می‌کند و برای تصویر $800 \times 800$
به ۵ دقیقه می‌رسد. بنابراین، فراتر از این اندازه، ZSSR سریع‌تر از
EDSR+ عمل می‌کند.

\subsection{سازگاری با تصویر آزمون}
\label{subsec:adaptation}

یکی از مزایای بارز ZSSR نسبت به روش‌های نظارت‌شده، قابلیت سازگاری
با تخریب‌ها و تنظیمات خاص هر تصویر آزمون در زمان آزمون است. در
روش‌های نظارت‌شده، شبکه برای یک هستهٔ کاهش‌نمونه‌برداری ثابت (معمولاً
دومکعبی) و شرایط تصویربرداری با کیفیت بالا بهینه شده است. در عمل،
فرایند تصویربرداری از تصویری به تصویر دیگر متفاوت است: دوربین‌ها و
حسگرها متنوع‌اند (لنزها و PSFهای مختلف)، شرایط فردی تصویربرداری
متغیر است (مثلاً لرزش ناخواستهٔ دوربین، شرایط دید ضعیف) و
این تفاوت‌ها به هسته‌های کاهش‌نمونه‌برداری مختلف، ویژگی‌های نویزی
گوناگون و آرتیفکت‌های فشرده‌سازی مختلف منجر می‌شود.

آموزش عملی برای تمام ترکیب‌های ممکن تنظیمات تصویربرداری غیرممکن
است. علاوه بر آن، یک CNN نظارت‌شدهٔ واحد بعید است برای تمام انواع
تخریب‌ها/تنظیمات عملکرد خوبی داشته باشد. دستیابی به عملکرد مطلوب
نیاز به شبکه‌های تخصصی متعدد دارد که هر یک بر نوع خاصی از تخریب‌ها
آموزش دیده باشند (برای روزها یا هفته‌ها). ZSSR با اجازه دادن به
کاربر برای تنظیم پارامترهای زیر در زمان آزمون، این مشکل را حل
می‌کند:

\begin{enumerate}
\item \textbf{هستهٔ کاهش‌نمونه‌برداری مطلوب:} هستهٔ مورد نظر
(در صورت عدم ارائه، هستهٔ دومکعبی به‌عنوان پیش‌فرض استفاده می‌شود).
\item \textbf{ضریب مقیاس ابرتفکیک $s$:} ضریب بزرگ‌نمایی مطلوب.
\item \textbf{تعداد افزایش‌های تدریجی مقیاس:} تعادل بین سرعت و
کیفیت (پیش‌فرض ۶).
\item \textbf{اعمال بازافکنی:} آیا قید سازگاری بین تصاویر LR و HR
اعمال شود (پیش‌فرض: بله).
\item \textbf{افزودن نویز به فرزندان LR:} آیا نویز مصنوعی به
فرزندان LR در هر جفت نمونهٔ آموزشی افزوده شود (پیش‌فرض: خیر).
\end{enumerate}

دو پارامتر آخر (حذف بازافکنی و افزودن نویز) امکان مدیریت ابرتفکیک
تصاویر LR با کیفیت پایین را فراهم می‌سازند. نویسندگان دریافتند که
افزودن مقدار کمی نویز گاوسی (با میانگین صفر و انحراف استاندارد
$\sim 5$ سطح خاکستری) عملکرد را برای طیف وسیعی از تخریب‌ها بهبود
می‌بخشد. علت این پدیده آن است که اطلاعات خاص تصویر تمایل به تکرار
در مقیاس‌های مختلف دارند، در حالی که آرتیفکت‌های نویزی این ویژگی
را ندارند \cite{zontak2013separating}. افزودن نویز مصنوعی به
فرزندان LR (بدون افزودن به پدرهای HR) به شبکه می‌آموزد که
اطلاعات نامرتبط بین‌مقیاسی (نویز) را نادیده بگیرد و تنها وضوح
اطلاعات مرتبط (جزئیات سیگنال) را افزایش دهد.

شایان ذکر است که ارائهٔ هستهٔ کاهش‌نمونه‌برداری تخمین‌زده‌شده به
روش‌های نظارت‌شدهٔ پیشرفته در زمان آزمون سودمند نخواهد بود، زیرا
آن‌ها نیاز به بازآموزش کامل شبکه بر مجموعهٔ جدیدی از جفت‌های
LR-HR دارند که با این هستهٔ خاص (ناپارامتری) تولید شده باشند.

\subsection{آزمایش‌ها و نتایج}
\label{subsec:zssr_experiments}

روش ZSSR عمدتاً برای تصاویر LR واقعی با تنظیمات تصویربرداری ناشناخته
و متغیر طراحی شده است. این تصاویر معمولاً فاقد حقیقت زمینی HR هستند
و ارزیابی بصری روش اصلی سنجش آن‌هاست. با این حال، برای ارزیابی کمّی،
چندین آزمایش کنترل‌شده در تنظیمات مختلف انجام شده است.

\subsubsection{حالت ایده‌آل}

هرچند هدف اصلی ZSSR حالت ایده‌آل نیست، عملکرد آن بر مجموعه‌داده‌های
استاندارد ابرتفکیک ایده‌آل نیز ارزیابی شده است. در این
مجموعه‌داده‌ها، تصاویر LR با دستور \texttt{imresize} متلب (هستهٔ
دومکعبی با ضدتداخل) از نسخهٔ HR تولید شده‌اند.

\begin{table}[!ht]
\centering
\caption{مقایسهٔ نتایج ابرتفکیک در حالت ایده‌آل (کاهش‌نمونه‌برداری
دومکعبی) بر حسب \lr{PSNR / SSIM}. برگرفته از \cite{shocher2018zssr}.}
\label{tbl:zssr_ideal}
\renewcommand{\arraystretch}{1.4}
\begin{tabular}{|c|c|c|c|c|c|}
\hline
\textbf{مجموعه‌داده} & \textbf{مقیاس} &
\textbf{SRCNN} & \textbf{VDSR} & \textbf{EDSR+} &
\textbf{ZSSR} \\
\hline
\multirow{3}{*}{\lr{Set5}}
& $\times 2$ & \lr{36.66/0.954} & \lr{37.53/0.959} &
\lr{\textbf{38.20}/\textbf{0.961}} & \lr{37.37/0.957} \\
& $\times 3$ & \lr{32.75/0.909} & \lr{33.66/0.921} &
\lr{\textbf{34.76}/\textbf{0.929}} & \lr{33.42/0.919} \\
& $\times 4$ & \lr{30.48/0.863} & \lr{31.35/0.884} &
\lr{\textbf{32.62}/\textbf{0.898}} & \lr{31.13/0.880} \\
\hline
\multirow{3}{*}{\lr{Set14}}
& $\times 2$ & \lr{32.42/0.906} & \lr{33.03/0.912} &
\lr{\textbf{34.02}/\textbf{0.920}} & \lr{33.00/0.911} \\
& $\times 3$ & \lr{29.28/0.821} & \lr{29.77/0.831} &
\lr{\textbf{30.66}/\textbf{0.848}} & \lr{29.80/0.830} \\
& $\times 4$ & \lr{27.49/0.750} & \lr{28.01/0.767} &
\lr{\textbf{28.94}/\textbf{0.790}} & \lr{28.01/0.765} \\
\hline
\multirow{3}{*}{\lr{BSD100}}
& $\times 2$ & \lr{31.36/0.888} & \lr{31.90/0.896} &
\lr{\textbf{32.37}/\textbf{0.902}} & \lr{31.65/0.892} \\
& $\times 3$ & \lr{28.41/0.786} & \lr{28.82/0.798} &
\lr{\textbf{29.32}/\textbf{0.810}} & \lr{28.67/0.795} \\
& $\times 4$ & \lr{26.90/0.710} & \lr{27.29/0.725} &
\lr{\textbf{27.79}/\textbf{0.744}} & \lr{27.12/0.721} \\
\hline
\end{tabular}
\end{table}

نتایج جدول \ref{tbl:zssr_ideal} نشان می‌دهد که ZSSR خاصِ تصویر، در
مقابل روش‌های نظارت‌شده‌ای که به‌صورت گسترده برای همین شرایط آموزش
دیده‌اند، نتایج رقابتی کسب می‌کند. به‌طور خاص، ZSSR به‌مراتب بهتر از
SRCNN \cite{dong2014srcnn} عمل می‌کند و در برخی موارد با
VDSR \cite{kim2016vdsr} (که تا یک سال قبل پیشرفته‌ترین روش بود)
قابل مقایسه یا حتی بهتر است. در رژیم ابرتفکیک بدون نظارت نیز
ZSSR با اختلاف زیاد از روش پیشرو \lr{SelfExSR} \cite{huang2015selfex}
پیشی می‌گیرد.

نکتهٔ قابل توجه آن است که در تصاویری با ساختارهای تکرارشوندهٔ
درونی قوی، ZSSR حتی از VDSR و گاهی از EDSR+ نیز فراتر می‌رود،
حتی اگر تصاویر LR با تنظیمات ایده‌آل نظارت‌شده تولید شده باشند.
تحلیل بیشتر نشان می‌دهد که برخی پیکسل‌ها (به‌ویژه در نواحی با
تکرارپذیری بالای اطلاعات، مانند پنجره‌های بسیار کوچک در بالای
ساختمان) از یادگیری داخلی (ZSSR) بیشتر سود می‌برند، در حالی که
پیکسل‌های دیگر (مثلاً لبه‌هایی که در تصاویر طبیعی فراوان هستند)
از یادگیری خارجی (EDSR+) بهره‌مند‌ترند. این مشاهده حاکی از
پتانسیل ترکیب یادگیری داخلی و خارجی در یک چارچوب محاسباتی واحد
برای بهبود بیشتر عملکرد است.

\subsubsection{حالت غیرایده‌آل}

تصاویر LR واقعی معمولاً به‌صورت ایده‌آل تولید نمی‌شوند. ZSSR در دو
دستهٔ اصلی از شرایط غیرایده‌آل ارزیابی شده است:
(الف) هسته‌های کاهش‌نمونه‌برداری غیرایده‌آل و
(ب) تصاویر LR با کیفیت پایین.

\paragraph{هسته‌های کاهش‌نمونه‌برداری غیرایده‌آل.}
برای ایجاد یک مجموعه‌دادهٔ کنترل‌شده، از مجموعهٔ BSD100 \cite{martin2001bsd}
استفاده شد. تصاویر HR با هسته‌های گاوسی تصادفی (با اندازهٔ معقول)
کاهش‌نمونه‌برداری شدند. ماتریس کوواریانس $\boldsymbol{\Sigma}$ هستهٔ
کاهش‌نمونه‌برداری هر تصویر به‌صورت تصادفی و مطابق رابطهٔ زیر انتخاب شد:
\begin{align}
\lambda_1, \lambda_2 &\sim \mathcal{U}\!\left[0, \tfrac{s}{2}\right],
\qquad \theta \sim \mathcal{U}[0, \pi],
\nonumber\\[4pt]
\boldsymbol{\Lambda} &= \mathrm{diag}(\lambda_1, \lambda_2),
\qquad
\mathbf{U} =
\begin{bmatrix}
\cos\theta & -\sin\theta \\
\sin\theta & \cos\theta
\end{bmatrix},
\nonumber\\[4pt]
\boldsymbol{\Sigma} &= \mathbf{U}\,\boldsymbol{\Lambda}\,\mathbf{U}^{\!\top},
\label{eq:zssr_kernel}
\end{align}
که $s$ ضریب کاهش‌نمونه‌برداری HR-LR است. بنابراین هر تصویر LR با یک
هستهٔ تصادفی متفاوت کاهش‌نمونه‌برداری شده است.

\begin{table}[!ht]
\centering
\caption{ابرتفکیک در حضور هسته‌های کاهش‌نمونه‌برداری ناشناخته بر
مجموعهٔ BSD100 ($\times 2$). برگرفته از \cite{shocher2018zssr}.}
\label{tbl:zssr_kernel}
\renewcommand{\arraystretch}{1.4}
\begin{tabular}{|c|c|c|}
\hline
\textbf{روش} & \textbf{\lr{PSNR}} & \textbf{\lr{SSIM}} \\
\hline
\lr{VDSR} \cite{kim2016vdsr} & \lr{27.72} & \lr{0.764} \\
\lr{EDSR+} \cite{lim2017edsr} & \lr{27.78} & \lr{0.766} \\
\lr{Blind-SR} \cite{michaeli2013blind} & \lr{28.42} & \lr{0.783} \\
\lr{ZSSR [هستهٔ تخمینی]} & \lr{28.81} & \lr{0.831} \\
\lr{ZSSR [هستهٔ واقعی]} & \lr{\textbf{29.68}} & \lr{\textbf{0.841}} \\
\hline
\end{tabular}
\end{table}

نتایج جدول \ref{tbl:zssr_kernel} حاکی از آن است که ZSSR با اختلاف
بسیار زیاد از روش‌های پیشرفته پیشی می‌گیرد:
$+1$~dB برای هسته‌های ناشناخته (تخمین‌زده‌شده با \cite{michaeli2013blind})
و $+2$~dB هنگام ارائهٔ هسته‌های واقعی. از دیدگاه بصری نیز تصاویر SR
تولیدشده توسط روش‌های نظارت‌شده بسیار تار هستند، در حالی که ZSSR
جزئیات واضح‌تری بازسازی می‌کند.

نکتهٔ جالب آن است که روش بدون نظارت \lr{Blind-SR} \cite{michaeli2013blind}
که از یادگیری عمیق استفاده نمی‌کند نیز از روش‌های نظارت‌شدهٔ
پیشرفته بهتر عمل می‌کند. این یافته از تحلیل‌ها و مشاهدات
\cite{efrat2013accurate} پشتیبانی می‌کند مبنی بر اینکه
(الف) مدل دقیق کاهش‌نمونه‌برداری مهم‌تر از پیشین‌های پیچیدهٔ تصویر
است و (ب) استفاده از هستهٔ نادرست منجر به نتایج بیش‌هموار‌شده می‌گردد.

حالت خاص هستهٔ $\delta$ (منجر به تداخل‌گذاری\LTRfootnote{Aliasing})
نیز توسط روش‌های نظارت‌شده به‌خوبی مدیریت نمی‌شود، در حالی که ZSSR
قادر به مدیریت آن است.

\paragraph{تصاویر LR با کیفیت پایین.}
در این آزمایش، برای هر تصویر از مجموعهٔ BSD100 \cite{martin2001bsd}
یک نوع تخریب تصادفی از میان سه نوع انتخاب شد:
(i) نویز گاوسی با $\sigma = 0.05$،
(ii) نویز لکه‌ای\LTRfootnote{Speckle noise} با $\sigma = 0.05$،
(iii) فشرده‌سازی \lr{JPEG} با کیفیت ۴۵.

\begin{table}[!ht]
\centering
\caption{ابرتفکیک در حضور تخریب ناشناختهٔ تصویر بر مجموعهٔ BSD100
($\times 2$). برگرفته از \cite{shocher2018zssr}.}
\label{tbl:zssr_degradation}
\renewcommand{\arraystretch}{1.4}
\begin{tabular}{|c|c|c|}
\hline
\textbf{روش} & \textbf{\lr{PSNR}} & \textbf{\lr{SSIM}} \\
\hline
درون‌یابی دومکعبی & \lr{27.92} & \lr{0.750} \\
\lr{EDSR+} \cite{lim2017edsr} & \lr{27.56} & \lr{0.714} \\
\lr{ZSSR} & \lr{\textbf{28.61}} & \lr{\textbf{0.781}} \\
\hline
\end{tabular}
\end{table}

نتایج جدول \ref{tbl:zssr_degradation} نشان می‌دهد که ZSSR در مقابل
تخریب‌های ناشناخته مقاوم است، در حالی که روش‌های نظارت‌شدهٔ پیشرفته
به حدی آسیب می‌بینند که حتی درون‌یابی دومکعبی ساده از آن‌ها بهتر عمل
می‌کند. این یافته به‌وضوح شکنندگی روش‌های نظارت‌شده در شرایط واقعی
را نمایان می‌سازد.

\subsection{تحلیل مزایا و محدودیت‌ها}
\label{subsec:zssr_analysis}

مزایای کلیدی روش ZSSR را می‌توان به‌صورت زیر خلاصه کرد:

\begin{enumerate}
\item \textbf{نخستین روش ابرتفکیک مبتنی بر CNN بدون نظارت:}
ZSSR نشان داد که امکان بهره‌گیری از قدرت یادگیری عمیق بدون نیاز
به هیچ دادهٔ آموزشی خارجی وجود دارد.
\item \textbf{مدیریت شرایط غیرایده‌آل:} این روش قادر به مدیریت
طیف وسیعی از شرایط تصویربرداری غیرایده‌آل از جمله هسته‌های
کاهش‌نمونه‌برداری ناشناخته، نویز حسگر، فشرده‌سازی و تداخل‌گذاری
است.
\item \textbf{عدم نیاز به پیش‌آموزش:} شبکه در زمان آزمون آموزش
می‌بیند و به منابع محاسباتی متوسط نیاز دارد.
\item \textbf{انعطاف‌پذیری مقیاس:} قابل اعمال برای هر اندازه و
نظریتاً هر نسبت ابعاد.
\item \textbf{سازگاری در زمان آزمون:} قابلیت تنظیم پارامترها
(هسته، نویز، بازافکنی) بر اساس ویژگی‌های تصویر آزمون.
\item \textbf{عملکرد پیشرفته در شرایط غیرایده‌آل:} بهبود
$1$--$2$~dB نسبت به روش‌های نظارت‌شدهٔ پیشرفته بر تصاویر غیرایده‌آل.
\end{enumerate}

در عین حال، ZSSR محدودیت‌هایی نیز دارد. در حالت ایده‌آل (هستهٔ
دومکعبی، بدون نویز)، عملکرد آن از روش‌های نظارت‌شدهٔ بسیار عمیق مانند
EDSR+ فاصله دارد، هرچند نتایج رقابتی هستند. علاوه بر آن، آموزش در
زمان آزمون هزینهٔ زمانی اضافی دارد، هرچند همان‌طور که بحث شد، این
هزینه برای تصاویر بزرگ قابل مقایسه یا حتی کمتر از زمان استنتاج
روش‌های بسیار عمیق است.

\subsection{نتیجه‌گیری و تأثیر بر پژوهش‌های بعدی}
\label{subsec:zssr_conclusion}

روش ZSSR \cite{shocher2018zssr} مفهوم ابرتفکیک «بدون نمونهٔ آموزشی»
را معرفی کرد که از قدرت یادگیری عمیق بدون اتکا به هیچ نمونهٔ خارجی
یا پیش‌آموزش بهره می‌گیرد. این امر از طریق یک شبکهٔ CNN کوچک خاص
تصویر محقق شد که در زمان آزمون بر نمونه‌های درونی استخراج‌شده صرفاً
از تصویر LR آزمون آموزش می‌بیند. این رویکرد ابرتفکیک تصاویر واقعی
را ممکن ساخت؛ تصاویری که فرایند تصویربرداری آن‌ها غیرایده‌آل، ناشناخته
و از تصویری به تصویر دیگر متفاوت است. در این شرایط غیرایده‌آل واقعی،
ZSSR هم از نظر کیفی و هم از نظر کمّی به‌طور قابل توجهی از روش‌های
پیشرفتهٔ ابرتفکیک پیشی گرفت.

مهم‌تر از نتایج کمّی، ZSSR پارادایم جدیدی را در ابرتفکیک تصویر
بنیان‌گذاری کرد: ترکیب قدرت تعمیم‌پذیری یادگیری عمیق با بهره‌گیری از
آمار درونی تصویر. این پارادایم مسیر را برای پژوهش‌های بعدی هموار
ساخت، از جمله روش‌هایی که یادگیری داخلی و خارجی را در یک چارچوب واحد
ترکیب می‌کنند و شبکه‌هایی که در زمان آزمون خود را با شرایط خاص
تصویر سازگار می‌سازند. این پژوهش همچنین اهمیت مدل دقیق فرایند
تخریب تصویر را بیش از پیشین‌های پیچیدهٔ تصویر برجسته ساخت؛ مشاهده‌ای
که تأثیر عمیقی بر جهت‌گیری تحقیقات ابرتفکیک در سال‌های بعد داشت.
